\section[Lagrange 量子场论]{Lagrange 量子场论（2 月 9 日）}

\srcpage{21}

今天我们从 Lagrange 的视角讨论欧氏量子场论（Euclidean quantum field theory，QFT）。

固定一个 $d$ 维流形 $M$；这里 $d$ 是时空的维数。我们把一个场空间 $\mathcal{F}(M)$ 附着到 $M$ 上。场是那些可以层化的东西，即可以由能粘接的局部数据拼装起来：例如主丛与联络、函数、向量丛的截面、微分形式等等。与场 $\varphi\in\mathcal{F}(M)$ 相伴有一个作用量（action）$S(\varphi)$，即函数 $S\colon\mathcal{F}(M)\to\R$。

量子场论中的「量子」意味着：我们不是计算确定的测量值，而是得到一个概率分布，办法是在场空间上积分。这就是臭名昭著的路径积分；它并不总有良好的数学定义，但眼下我们不管这个。物理学家有计算路径积分的办法。

若 $M$ 闭，我们能做的最基本的测量是场空间的体积，即配分函数
\[
  Z(M) := \int_{\mathcal{F}(M)} e^{-S(\varphi)/\hbar}\,\mathcal{D}\varphi.
\]

我们还想计算算子的期望值。若 $\mathcal{O}_x$ 是点 $x\in M$ 处的算子，它定义了 $\mathcal{F}(M)$ 上的一个泛函，物理上可以解释为在 $x$ 处做一次测量。该测量由修改过的路径积分算出：
\[
  Z(M) := \int_{\mathcal{F}(M)} \mathcal{O}_x(\varphi)\, e^{-S(\varphi)/\hbar}\,\mathcal{D}\varphi.
\]

其中有无序算子（disorder operator）：插入 $\mathcal{O}_x$ 意味着考虑可能在 $x$ 处奇异的场，例如考虑 $\mathcal{F}(M\setminus x)$ 中的好场。

\srcpage{22}

我们不打算把它严格化，但这有助于理解时间演化以及上次的 Hamilton 视角。设 $M$ 是从 $\partial M_{\mathrm{in}}$ 到 $\partial M_{\mathrm{out}}$ 的协边。因为场可以层化，可以把它们限制到入、出边界，得到映射 $\mathcal{F}(\partial M_{\mathrm{in}})\leftarrow\mathcal{F}(M)\to\mathcal{F}(\partial M_{\mathrm{out}})$。我们又（或多或少）有 $\mathcal{F}(M)$ 上的路径积分测度 $\mathcal{D}\varphi$，于是可以做积分变换——至少在启发式层面——从 $\mathcal{F}(\partial M_{\mathrm{in}})$ 上的泛函变到 $\mathcal{F}(\partial M_{\mathrm{out}})$ 上的泛函。记这两个泛函空间为 $\mathcal{H}_{\mathrm{in}}$ 与 $\mathcal{H}_{\mathrm{out}}$。这里要操心的细节很多，包括正则性以及路径积分测度的定义。

无论如何，从这些数据我们得到一个线性映射 $\mathcal{H}_{\mathrm{in}}\to\mathcal{H}_{\mathrm{out}}$。它在泛函 $f\in\mathcal{H}_{\mathrm{in}}$ 与场 $\varphi\in\mathcal{F}(\partial M_{\mathrm{out}})$ 上的值可以这样描述：考虑 $\mathcal{F}(M)$ 中那些连同在 $\partial M_{\mathrm{out}}$ 上限制为 $\varphi$ 的映射一起的场构成的子空间 $T$；然后用 $\partial M_{\mathrm{in}}$ 的路径积分测度，对 $T$ 中诸场限制到 $\partial M_{\mathrm{in}}$ 上、$f$ 的取值作积分。

这就是 Lagrange 视角与 Hamilton 视角之间的关系。若 $M$ 是解释为恒等协边的柱体，我们希望把这个映射看成时间演化 $e^{itH/\hbar}$，其中 $t$ 是柱体的长度。Lagrange 视角揭示出许多额外的丰富结构，我们会发现它们很有用。

\begin{example}
  设 $d = 1$，我们想研究 Riemann 流形 $X$ 上的量子力学。那么 $1$-流形 $C$ 上的场是 $\mathcal{F}(C) := \mathrm{Map}(C,X)$；作用量的临界点是测地线；$\mathcal{H} = L^2(\mathrm{Map}(\mathrm{pt},X)) = L^2(X)$，而 Hamilton 量是 Laplacian 算子。
\end{example}

然而，若想在更高维做这件事，由无限维场空间带来的问题会更严重。但我们还是继续往前走，看看即使不数学精确也能学到什么。

\begin{example}[量子 Maxwell 理论]
  这是一个四维场论。$M$ 上的场是 $M$ 上带联络 $d + A$ 的主 $U_1$-丛（的规范等价类）。记 $F$ 为联络的曲率。经典方程（即解为作用量临界点的方程）是 $dF = 0$ 与 $d\star F = 0$。作用量为
  \[
    S = \frac{g}{2\pi i}\int_M F\wedge\star F + \theta\int F\wedge F,
  \]
  其中 $g$ 与 $\theta$ 是参数，可能还有些常数要定。注意由 Chern--Weil 理论，$F\wedge F$ 是 $c_1(P)^2$ 的固定倍数，其中 $P\to M$ 是主 $U_1$-丛。

  在闭 $3$-流形 $N$ 上，Hilbert 空间 $\mathcal{H}$ 是场空间 $\mathcal{C}(N)$ 上的函数，具体来说是主 $U_1$-丛上联络空间模去规范群后的 $L^2$。

  因为 $U_1$ 阿贝尔，我们想把 $\mathcal{C}(N)$ 看成某种无限维 Lie 群。取第一 Chern 类给出满射 $\mathcal{C}(N)\twoheadrightarrow \Lambda := H^2(N;\Z)$，\yizhu{原文此处写作 $H^2(N:\Z)$，用的是冒号而非分号，应为笔误；译文按分号处理。}{A（原文 p.22）}这个映射的核是平坦联络的空间 $T := H^1(N;U_1)\cong BH^1(N;\Z)$。场空间分裂为 $\Lambda\times T\times V$，其中 $V$ 是某个无限维向量空间。

  格 $\Lambda$ 作用在 $\mathcal{H}$ 上，这与 Dirac 单极子有关；不过在非阿贝尔情形，这与 't Hooft 算子有关。具体地说，给定 $\lambda\in\Lambda$，有线丛 $L$ 使 $c_1(P) = \lambda$。于是 $\Lambda$ 对 $\mathcal{C}(N)$ 的作用就是与 $L$ 取张量积，这诱导了 $\mathcal{H}$ 上的作用。

  这在 $M$ 中造出磁单极子！设 $\gamma$ 是光滑闭曲线，其同调类 Poincaré 对偶于 $\lambda$。（TODO：物理内容该放这里——这与 $\gamma$ 上的磁单极子有关。）

  现在若取 $N$ 的柱体协边 $M\times I$，并在 $t = 1/2$ 处嵌入一个结 $\gamma$\yizhu{原文此句写作「if $N$ is the cylinder bordism for $N$」，主语重复，按上下文应为 $M\times I$ 是 $N$ 的柱体协边；译文按上下文处理。}{A（原文 p.22）}，我们的方程排除了磁单极子，所以需要考虑 $(M\times I)\setminus\gamma$ 上的场空间。这个结的链环是 $S^2\times S^1$。在 $\gamma$ 上取我们「准单极子」所在的一点，令 $S$ 为环绕 $2$-球面。则 $S$ 内部的磁荷非零。

  来看联络构成的空间 $\mathcal{C}_{c=1}$：它是 $(M\times I)\setminus\gamma$ 上满足 $\frac{1}{2\pi i}\int_{S^2} F = 1$ 的联络，即 $2$-球面内部有一个单位的磁荷。这个联络空间在数学上有良好定义，而我们得到一个对应
  \[
    \mathcal{C}(N_{\mathrm{in}})\longleftarrow\mathcal{C}_{c=1}\longrightarrow\mathcal{C}(N_{\mathrm{out}}).
  \]

  因此由通常的积分变换，得到映射 $\mathcal{H}\to\mathcal{H}$，其效果是平移 $\lambda = c_1(P)$。这是无序算子的一个例子：我们不是在测量场，而是在看另一个场空间。我们的测量使场沿 $\gamma$ 变奇异。场论式的描述很复杂，尽管它在 $\mathcal{H}$ 上的作用很简单。
\end{example}

\srcpage{23}

Wilson 算子的场论描述相应地简单。设 $\gamma\subset N$ 是简单闭曲线。即使在经典层面，也有函数 $W_\gamma\colon\mathcal{C}(N)\to\C$，把 $(L,\nabla)$ 送到 $\nabla$ 绕 $\gamma$ 的和乐。它取值在 $U_1\subset\C$ 中。我们同样可以把 $\gamma$ 放在 $[0,1]\times M$ 中 $t = 1/2$ 处，从而定义积分变换 $\mathcal{H}\to\mathcal{H}$。我们不再是在一个点处测量，而是沿一个环路测量。其作用是让 $\mathcal{H}$ 中的泛函乘以 $W_\gamma$。

这些 $W_\gamma$ 是 $T$ 在 $L^2(\mathcal{C}(M))$ 上作用（乘以和乐）的特征函数。这是该验证一下的，但其实就是 Fourier 理论。它给出分解
\[
  \mathcal{H} = \bigoplus_{e\in H(M;\Z)} \mathcal{H}_e,
\]
即电通量。

这是电磁对偶的第一部分；可以把它看成在 $\mathcal{C}(N)$ 上做 Fourier 分析，而我们已经把它分解成格、环面与向量空间。Fourier 变换把 $L^2(\mathcal{C}_{U_1}(M))$ 与 $L^2(\mathcal{C}_{U_1^\vee}(M))$ 等同，把格 $\Lambda_B$ 与环面 $T_B^\vee$ 互换，反过来把 $T_E$ 与 $\Lambda_E^\vee$ 互换，而向量空间映到向量空间。

在一侧，Wilson 算子是对角的；在另一侧，'t Hooft 算子是对角的。这是电磁对偶的标志。

电磁对偶本质上是量子的：你无法从经典的 Lagrange 视角看到它。必须量子化才能看见其中一半的算子。从 Lagrange 视角看，磁算子通过让联络沿环路奇异而破坏它。既然电磁对偶把它换成经典视角下说得通的算子，我们就需要另一个视角才能看见它。

后面讲到几何 Satake 同构时，我们会把它解释为：告诉我们非阿贝尔规范场论中 't Hooft 算子是什么。

\begin{example}[$d = 2$ 中的阿贝尔对偶]
  这将是 T-对偶的一个简单实例。设 $M$ 是 $2$-流形；场是 $\mathrm{Map}(M,S^1)$，而 $S^1$ 上的 Hilbert 空间是 $L^2(\mathrm{Map}(S^1,S^1))$。它由 $H^1(S^1;\Z)$ 分次，即映射的绕数。

  T-对偶在这里的作用是：可以取对偶圆 $(S^1)^\vee$；若原圆的半径是 $R$，对偶圆半径是 $1/R$。于是这两个映射空间的 $L^2$ 之间有对应。具体地，对偶情形可以考虑 $L^2(H^0(M;U_1))$。

  在原侧，我们可以平移绕数，这是 $4$ 维 Maxwell 理论中 't Hooft 算子的类比。在对偶侧，我们可以在一点处求值，这是 Wilson 的类比。对偶互换这些算子，是镜像对称的极简情形。
\end{example}

本课最关心的量子场论是拓扑的。我们的目标是从无限维流形上函数构成的巨大 Hilbert 空间，砍到有限维 Hilbert 空间，而且要以仍然有用的方式砍。

例如，看量子力学。点被赋予一个 Hilbert 空间 $\mathcal{H}$，时间区间 $[0,T]$ 被赋予时间演化算子 $e^{iT\mathcal{H}/\hbar}$。Hamilton 量是 Laplacian 算子，Laplacian 的整个谱理论都出现在故事里，其中有微妙之处。

在拓扑量子力学中，我们把时间杀死。不是「我们多待一会儿」那个意思，而是「我们强迫时间演化成为恒等映射」。有几种做法：可以限制到 Laplacian 的零特征子空间，但那些是调和函数；由极大值原理，它们在闭 Riemann 流形上是常数。于是得到的理论并不特别有趣。另一种做法是照 Witten [Wit82]，通过引入超对称，在导出意义下杀死时间。

（1）首先，把 $L^2(X)$ 扩大为目标 $X$ 上 $L^2$ 微分形式的空间。于是我们加了新的场。（2）利用这一点，可以找出对称性的（超）Lie 群。

对 $\Omega^\bullet_X$，除了按分次作用的 $\mathfrak{u}_1$ 之外，还有两个算子 $Q := d$ 与 $Q^* := d^*$。$Q$、$Q^*$ 与 $\mathfrak{u}_1$（R-对称性）之间的交换关系告诉我们，应把 $Q$ 与 $Q^*$ 分别放在次数 $1$ 与 $-1$。它们共同定义一个分次 Lie 代数，称为一维 $\mathcal{N} = 1$ 超对称代数；它们彼此分次交换，只有 $[Q,Q^*] = H$：$\Delta = dd^* + d^*d$。其余的交换关系用到 $d^2 = 0$ 等等。

也就是说，$Q^*$ 是从 $H$ 到 $0$ 的（链）同伦。因此在 $Q$-上同调（即通常的 de Rham 上同调）上，$H$ 作用为零：这里有点 Hodge 理论的味道，说调和形式在 de Rham 上同调中给出代表元。
\srcpage{24}
然而我们还剩下更多东西：也许不是调和函数，但现在我们有上同调！

也就是说，拓扑量子力学把「Hilbert 空间」$(\Omega^\bullet_X, d)$，或 $H^*_{dR}(X)$，赋予 $X$，此时 $H = 0$。而这是 $X$ 的拓扑不变量，表明对 Riemann 几何的依赖已经消失。这非常激进！
